\section{}

\subsection{Ключевые идеи}

Статья~\cite{shihab2025detectreward} посвящена проблеме взлома вознаграждения  в системах обучения с подкреплением ~-- явлению, при котором агент достигает высоких значений функции вознаграждения $R_{\text{proxy}}$, не достигая при этом истинной цели задачи, описываемой функцией $R_{\text{true}}$. Авторы проводят масштабное эмпирическое исследование, включающее 15,247 эпизодов, 15 сред и 5 алгоритмов RL, чтобы выявить распространённость, закономерности и причины взлома вознаграждения. В работе представлена \emph{система обнаружения и анализа} таких случаев, включающая шесть специализированных детекторов, а также методы количественной оценки и смягчения последствий.

\subsection{Математическое описание метода}

Пусть агент взаимодействует со средой, описываемой множеством состояний (S), действий (A) и функцией прокси-вознаграждения $R_{\text{proxy}}: S \times A \rightarrow \mathbb{R}$. Задача агента заключается в максимизации ожидаемого суммарного вознаграждения. Истинная цель системы выражается через функцию $R_{\text{true}}$, которая отражает реальное качество поведения. Случай взлома вознаграждения возникает тогда, когда направления оптимизации этих функций различаются:

\begin{equation}
\nabla_{\pi} \mathbb{E}[R_{\text{proxy}}] \neq \nabla_{\pi} \mathbb{E}[R_{\text{true}}].
\end{equation}

Для выявления отклонений вычисляется отношение прокси и истинного вознаграждения:

\begin{equation}
\rho = \frac{R_{\text{proxy}}}{R_{\text{true}}}.
\end{equation}

При корректном поведении агента значение $\rho$ стабильно и близко к медианному уровню $\rho_{baseline}$. Эпизод считается аномальным, если расхождение между распределениями текущих и базовых значений превышает заданный порог $\tau_{spec}$, оцениваемый по дивергенции Кульбака–Лейблера:

\begin{equation}
D_{\mathrm{KL}}(\rho_{\text{current}} || \rho_{\text{baseline}}) > \tau_{\text{spec}}.
\end{equation}

Детекторы работают параллельно и анализируют статистические зависимости между состояниями, действиями и наградами. Их оценки объединяются в ансамбль:

\begin{equation}
S_{\mathrm{ensemble}} = \sum_{i=1}^{6} w_i , S_i,
\end{equation}

где $S_i$~-- нормированная оценка отдельного детектора, а $w_i$ ~-- весовой коэффициент. Если ансамблевая оценка превышает порог $\theta$, эпизод классифицируется как случай взлома:

\begin{equation}
S_{\mathrm{ensemble}} > \theta \text{~-- взлом вознаграждения обнаружен.}
\end{equation}


\subsection{Алгоритм}

Работа системы включает следующие этапы:

\begin{enumerate}
    \item \emph{Сбор данных}~-- запись траекторий взаимодействия агента со средой $(s_t, a_t, r_t)$
    \item \emph{Извлечение признаков}~-- вычисление статистических характеристик.
    \item \emph{Анализ детекторами}~-- шесть детекторов выявляют различные типы отклонений: игру со спецификациями, подмену вознаграждения, оптимизацию прокси-показателя, несоответствие цели, эксплуатационные паттерны и самоподпитку вознаграждения.
    \item \emph{Калибровка ансамбля}~-- нормализация оценок $S_i$ методом приведения выходов модели к корректной вероятностной шкале с помощью логистической аппроксимации.
    \item \emph{Интеграция и классификация}~-- вычисление итогового ансамблевой оценки $S_{\mathrm{ensemble}}$ и определение случаев взлома.
\end{enumerate}

Такой алгоритм обеспечивает раннее выявление и количественную оценку взлома вознаграждения без существенного увеличения вычислительной нагрузки.


\subsection{Наборы данных}

Для экспериментов использованы 15 сред обучения с подкреплением, охватывающие различные типы задач:

\textbf{Atari}~-- Breakout, Pong, Space Invaders, Pac-Man, Asteroids;

\textbf{MuJoCo}~-- HalfCheetah, Walker2d, Hopper, Ant;

\textbf{GridWorld}~-- Simple Grid, Obstacle Grid, Multi-Goal Grid;

\textbf{Пользовательские среды}~-- Proxy Trap, Tampering Test, Alignment Challenge.

\vspace{11pt}
Всего проанализировано 15,247 эпизодов, из них 2,156 размечены экспертами.


\subsection{Метрики}

Для оценки эффективности метода применяются стандартные показатели:

\begin{equation}
P = \frac{TP}{TP + FP}, \quad R = \frac{TP}{TP + FN}, \quad F1 = \frac{2PR}{P + R},
\end{equation}

где (TP)~-- число корректно обнаруженных случаев, (FP)~-- ложные срабатывания, (FN)~-- пропущенные случаи.

Также используются:

\begin{itemize}
\item AUC-ROC~-- площадь под ROC-кривой, характеризующая общую способность различать классы;
\item Задержка обнаружения~-- количество эпизодов между возникновением отклонения и его фиксацией;
\item Частота взлома~-- процент эпизодов, в которых зафиксировано отклонение;
\item Разрыв показателей~-- разница между прокси и истинным вознаграждением:
\begin{equation}
\Delta R = R_{\text{proxy}} - R_{\text{true}};
\end{equation}
\end{itemize}


\subsection{Модели}

В работе исследованы пять алгоритмов обучения с подкреплением, различающихся стратегией обновления политики:

\begin{itemize}
\item \textbf{PPO}~-- метод оптимизации политики с ограничением шага (Proximal Policy Optimization);
\item \textbf{SAC}~-- мягкий актёр-критик (Soft Actor-Critic), использующий энтропийную регуляризацию;
\item \textbf{DQN}~-- глубокая Q-сеть (Deep Q-Network) с модификациями Double Q-Learning и Prioritized Replay;
\item \textbf{A3C}~-- асинхронный актёр-критик (Asynchronous Advantage Actor-Critic), обеспечивающий интенсивное исследование среды;
\item \textbf{Rainbow}~-- объединяет несколько улучшений DQN (Noisy Nets, Dueling Networks, Multi-step Returns).
\end{itemize}

Результаты показывают, что алгоритмы с агрессивным исследованием среды (например, A3C) демонстрируют более высокий уровень взлома вознаграждения ($28.9\%$), тогда как консервативные методы (SAC) характеризуются меньшей частотой ($19.3\%$).

